\section{Reducers checkpoints and replay}\label{sec:C13:replay}
Define a reducer $F_v$ that constructs application state from a prior state and an event under schema version $v$:
\begin{equation}
 S_t=F_v(S_{t-1},e_t),\qquad
 S_t=F_v^{*}(S_0,e_{1:t}).
\end{equation}
If $F_v$ is deterministic and all inputs are retained, replay can reconstruct the same application state. The assumptions matter. A reducer that silently consults the current clock, a mutable external service, or a newly upgraded language model is not deterministic under the recorded event stream. Nondeterministic results must be materialized as events if replay is expected to reproduce their consequences.

A checkpoint is a persisted state snapshot at a known event boundary. It can reduce replay cost from processing the complete history to processing the suffix after that boundary. Its identity should include the reducer version, event position or causal frontier, relevant schema versions, and references to external artifacts. A checkpoint that stores only the visible conversation can omit pending tool operations, unpublished artifacts, and reservations on a resource budget. Restoring that conversation may reproduce the model input while failing to restore the actual task.

There are at least three replay objectives. State replay reconstructs the recorded application state. Decision replay reconstructs the exact model input used for a past action and, where possible, its sampled output. Environment replay reexecutes actions against an equivalent environment snapshot. State replay does not imply decision replay if tokenization or context assembly changed. Decision replay does not imply environment replay if tools read changing data or produce irreversible effects. Evaluation should identify which objective it achieved.

A proposed checkpoint protocol writes an event frontier, state snapshot, artifact manifest, and integrity digest as one committed checkpoint object. An interrupted write should remain uncommitted rather than becoming an apparently valid partial snapshot. Restore verifies that all required objects exist and satisfy their recorded schemas. Missing optional caches can be recomputed; missing effect receipts or authority records may require stopping or external reconciliation. The distinction should be encoded instead of handled through a generic best-effort restore.

Checkpoint frequency trades storage and synchronization work against recovery time. If failures occur with a workload-dependent rate, a more frequent checkpoint can reduce lost computation but increase steady-state overhead. Language-agent workloads complicate this tradeoff because expensive steps are not uniform: one model call may be cheap while one external operation changes the world. A checkpoint before a consequential effect can be more valuable than one after a fixed number of tokens. Recovery cost should therefore be measured in both computation and semantic work lost.

LangGraph distinguishes checkpoints at super-step boundaries from per-task pending writes. Its documented durability modes also differ: synchronous mode completes persistence before the next step, asynchronous mode overlaps it with execution, and exit mode delays persistence until graph exit. An in-memory saver does not survive process restart. Moreover, replay from an earlier checkpoint reexecutes subsequent nodes, including model calls and API requests~\cite{cumb2026langgraphcheckpoints}. These are distinct recovery contracts, not interchangeable settings. They do not make an external effect atomic with checkpoint storage. The application must state the selected backend, durability mode, replay frontier, and treatment of side effects.

A concrete distinction appears in ACM's inspected \texttt{HistoryManager}. Its live message list is separate from an append-only raw mirror; \path{compress_range} removes a half-open slice only from the live list. However, \texttt{snapshot} returns the live messages, and \path{load_snapshot} initializes both lists from that snapshot. A different constructor, \path{from_saved_history}, accepts a separate raw history~\cite{repo_acm2026}. Therefore, using the snapshot path alone after compression need not preserve the original raw archive. This is a source-level contract difference, not a measured failure in the paper's experiments. A durable wrapper should persist both views when its recovery guarantee includes original evidence. Testing only that the live conversation restores would miss the archive distinction.

